\section{Introduction and main results}
\label{intro:section}

An invariant can become a pure power after reduction even though its
state differential, divided integrally before reduction, retains
nonzero information. For the root-of-unity q-Painlev\'e I invariant,
that divided differential reduces to a Hasse power times the
small-order invariant differential. This residue describes ordinary
levels but vanishes on supersingular levels. We determine the missing
vertical information: the complete coefficient ideal at the first
cyclotomic layer and the full first jet at every higher layer.

The first-layer result is controlled by a specific obstruction. The
original small-order invariant is a section of an anticanonical line
bundle that does not lift across the first infinitesimal cyclotomic
thickening. An integer congruence for the original matrix trace turns
its Bockstein into the connecting class of a locally exact differential
on each complete supersingular fiber. This differential is nowhere
zero. It supplies the second state direction that is absent from the
residue formula and determines the full local ideal. Higher layers
retain a common factor at odd primes, whereas characteristic two has
a different height-two base and a mixed cubic relation.

\subsection{The original invariant and its model}

We use the q-Painlev\'e I coordinates and root-of-unity integral of
Joshi--Roffelsen \cite{JoshiRoffelsen2026}. The one-step rational map is
\[
 F_t(x,y)=\left(\frac{t}{x-y/s},\frac{sx}{y}\right),\qquad t'=st.
\]
All parameters called time below are units. The invariant is fixed by
the original matrix, in the auxiliary normalization $w=1$.
Put $b=x(y-1)$, $d_0=y-1$, and $J_1(t)=y-x+x/y-t/x$. Then
\begin{equation}
 A(z;t)=
 \begin{pmatrix}t-b&-x\\d_0(b-t)&b\end{pmatrix}
 +z\begin{pmatrix}J_1(t)-1&1\\J_1(t)+b-1&1\end{pmatrix}
 +z^2\begin{pmatrix}1&0\\0&0\end{pmatrix}.
 \label{main:matrix}
\end{equation}
Its determinant is $z^3$. If $s$ has exact order $r$, the descending
product and the integral are
\begin{equation}
 M_{r,s}(z)=A(s^{r-1}z;t)\cdots A(z;t),\qquad
 I_{r,s}=[z^r]\operatorname{tr}M_{r,s}(z).
 \label{main:integral}
\end{equation}
The original identities are
\begin{equation}
 \operatorname{tr}M_{r,s}=t^r+I_{r,s}z^r+z^{2r},\qquad
 \det M_{r,s}=(-1)^{r+1}z^{3r},\qquad
 I_{r,s}=-t^r x^{-r}+O(x^{-r+1}).
 \label{main:original-identities}
\end{equation}
Here the last expression is a Laurent valuation statement, not an
analytic asymptotic. These matrix and integral identities are source
inputs, not new results of this paper; the formula numbering of the
author version specified in the bibliography is used when needed.

Let $S$ be the original eight-center blowup of
$\mathbb P^1\times\mathbb P^1$, let $D=-K_S$ be its boundary cycle,
and put $\mathcal U=S\setminus D$ and $L_n=\mathcal O_S(nD)$.
Section~\ref{geom:section}
gives the centers over the integer parameter ring and their terminal
charts. The open surface contains the torus and four complete terminal
affine lines. None of the results below removes these lines. We retain
the original time, energy, and two relative state directions; replacing
the matrix by an isospectral one would not define the same problem.

\subsection{Cyclotomic parameters and the three results}

Fix a prime $p$, an integer $m\ge1$ prime to $p$, and $a\ge1$.
Let $\mathcal O_0$ be an unramified $p$-adic discrete valuation ring
containing a primitive $m$th root $\widetilde\eta$, with perfect
residue field $k$. Set
\[
 \mathcal O_a=\mathcal O_0[\zeta_{p^a}],\quad
 \pi_a=\zeta_{p^a}-1,\quad s_a=\widetilde\eta\zeta_{p^a},\quad
 t_a\in\mathcal O_a^\times,\quad N=p^a.
\]
On the same eight-center model with $q=s_a$ and time $t_a$, define
\begin{equation}
 \alpha_{mN}=p^{-a}d_{\mathrm{state}}I_{mN,s_a}.
 \label{main:alpha}
\end{equation}
The differential fixes the time and the base ring. Its division occurs
in characteristic zero; Proposition~\ref{trace:regular} proves its
regularity on the entire $\mathcal U_a$.
With $\eta=\bar s_a$, put
\[
 J=I_{m,\eta}(x,y;\bar t_a),\qquad T=\bar t_a^m,\qquad
 \varepsilon_m=(-1)^{m+1},\qquad \sigma=\frac{p^a-1}{p-1}.
\]
The polynomial relevant to the complete original fibers is
\begin{equation}
 H_p(T,h;\varepsilon)=
 \begin{cases}
 [Z^{p-1}]\bigl((T+hZ+Z^2)^2-4\varepsilon Z^3\bigr)^{(p-1)/2},&p\ne2,\\
 h,&p=2.
 \end{cases}
 \label{main:hasse}
\end{equation}
Write $H=H_p(T,J;\varepsilon_m)$. It is monic of degree $p-1$ in
the energy variable, but its roots are not assumed simple.
Sections~\ref{spec:section}--\ref{hasse:section} prove that a complete
smooth finite fiber $X=(J=h)$ of the original pencil is supersingular
exactly when $H_p(T,h;\varepsilon_m)=0$.

For a regular one-form $\gamma$ on a smooth relative surface,
$\mathfrak c(\gamma)$ denotes the ideal generated by both coefficients
in a local frame of the relative cotangent sheaf. An invertible change
of frame does not change this ideal. The notation $\gamma^{[2]}$
means reduction modulo $\pi_a^2$, whose parameter image we denote by
$\epsilon$. If that quotient has characteristic $p$ and $b>0$ with
$p\mid b$, the $b$th power of any lift of a residue function
is independent of the lift: changing it by $\epsilon u$ changes the
power by $b\epsilon u$ times the preceding power, which is zero.
We write this well-defined lift-power as $f^{\langle b\rangle}$;
the zero power is one.

\begin{theorem}[All first layers]
\label{main:first}
Let $a=1$, with $k$ perfect. At every point of every original complete
smooth finite fiber, including its terminal points,
\begin{equation}
 \mathfrak c(\alpha_{mp})=(\pi_1,\widetilde H),
 \label{main:first-ideal}
\end{equation}
for any local lift $\widetilde H$ of $H$. If $h$ is a Hasse root
of multiplicity $e_h$, the completed ideal is $(\pi_1,z^{e_h})$,
where $\bar z=J-h$, after an allowed unramified residue-field extension.

On each complete smooth finite fiber, the actual Bockstein and the
restriction trivialized by the original section $1$ satisfy
\[
 0\ne\kappa_J=\rho_X\beta_{L_m}(J)\in H^1(X,\mathcal O_X).
\]
On a smooth supersingular fiber, the first tangent differential $\nu$
constructed from the original trace satisfies the exact comparison
\[
 \partial\nu=\operatorname{Fr}_*\kappa_J\ne0
 \quad\text{in }H^1(X,\mathcal O_X^p).
\]
The Frobenius here maps to its image sheaf, not to
$H^1(X,\mathcal O_X)$ after the subsequent inclusion.
\end{theorem}

The complete proof is in Section~\ref{first:section}. In particular,
\eqref{main:first-ideal} is an equality of the original full ideals,
not just of their images modulo $\pi_1^2$.

\begin{theorem}[Odd-prime higher first jets]
\label{main:odd}
Let $k$ be finite, $p$ odd, and $a\ge2$. Identify
$\mathcal O_1/(\pi_1^2)$ and $\mathcal O_a/(\pi_a^2)$ with
$k[\epsilon]/(\epsilon^2)$ by $\pi_1\mapsto\pi_a$, and choose
$t_1$ with the same full image as $t_a$ in this ring. On the entire
common open model,
\begin{equation}
 \alpha_{mp^a}^{[2]}=H^{\langle\sigma-1\rangle}\alpha_{mp}^{[2]}.
 \label{main:odd-form}
\end{equation}
Near every original complete smooth finite fiber,
\begin{equation}
 \mathfrak c(\alpha_{mp^a})+(\pi_a^2)
 =(\pi_a^2,\widetilde H^\sigma,\pi_a\widetilde H^{\sigma-1}).
 \label{main:odd-ideal}
\end{equation}
At a Hasse root of multiplicity $e_h$, its completed form is
$(\pi_a^2,z^{e_h\sigma},\pi_a z^{e_h(\sigma-1)})$.
\end{theorem}

The ring identification in this statement is not the natural embedding
of cyclotomic fields. Matching only the residue of the time is
insufficient; the complete time jet is matched. The form identity does
not require a smooth level or a nonzero Hasse value. The ideal assertion
uses smoothness and Theorem~\ref{main:first} after that identity has
been proved in Section~\ref{odd:section}.

\begin{theorem}[Characteristic-two higher first jets]
\label{main:two}
Let $p=2$, $m$ be odd, and $a\ge2$, with $N=2^a$.
Identify $\mathcal O_2/(\pi_2^2)$ and $\mathcal O_a/(\pi_a^2)$
with $B=k[\epsilon]/(\epsilon^2)$ by $\pi_2\mapsto\pi_a$,
and match the complete images of $t_2,t_a$. Over a perfect residue
field, on the entire common open model,
\begin{equation}
 \alpha_{mN}^{[2]}=J^{\langle N-4\rangle}\alpha_{4m}^{[2]}.
 \label{main:two-form}
\end{equation}
For finite $k$, near every original complete smooth finite fiber,
\begin{equation}
 \mathfrak c(\alpha_{mN})+(\pi_a^2)
 =(\pi_a^2,j^{N-1}+\pi_a\widetilde Tj^{N-4},\pi_a j^{N-2}),
 \label{main:two-ideal}
\end{equation}
where $j$ is any local lift of $J$ and $\widetilde T$ a constant
unit lift of $T$. At $a=2$ on $X=(J=0)$, the completed truncated
quotient, after the allowed unramified extension, is
\begin{equation}
 \frac{k(P)[[w,z,\epsilon]]}{(\epsilon^2,z^3+\epsilon T,\epsilon z^2)}
 \simeq\frac{k(P)[[w,z]]}{(z^5)},\qquad
 \pi_2\longmapsto z^3/T,
 \label{main:length-five}
\end{equation}
where $w$ is a parameter along the curve and $\bar z=J$.
\end{theorem}

The characteristic-two comparison starts at height two. The height-one
ring $\mathcal O_1/(\pi_1^2)$ has characteristic four and cannot be
used as its dual-number base. Section~\ref{two:section} computes the
actual four-block trace and its tangent class before eliminating the
two local coefficients. The integer five in \eqref{main:length-five}
is a transverse length of the truncated quotient: the parameter $w$
remains, and no length of the original full critical quotient is
asserted. After a finite unramified extension a cube root of $T$
normalizes its unit coefficient; $T$ is not a new formal modulus.

A state lift means a point over the same cyclotomic ring after a finite
unramified extension of its unramified coefficient ring. The common
order of a form is the minimum valuation of its two pulled-back state
coefficients, not the differential of the constant lifted point.
At ordinary smooth states, $\alpha_{mp^a}$ has order zero and
$dI_{mp^a}$ has order $a\varphi(p^a)$. At supersingular smooth states
the exact values and lower bounds are as follows.

\begin{table}[htbp]
 \centering
 \begin{tabular}{lll}
  \toprule
  Cyclotomic layer & Order of $\alpha_{mp^a}$ & Order of $dI_{mp^a}$\\
  \midrule
  Any $p$, $a=1$ & $1$ & $p$\\
  Odd $p$, $a\ge2$ & $\ge2$ & $\ge a\varphi(p^a)+2$\\
  $p=2$, $a=2$ & $1$ & $5$\\
  $p=2$, $a\ge3$ & $\ge2$ & $\ge a2^{a-1}+2$\\
  \bottomrule
 \end{tabular}
 \caption{Orders along the allowed unramified smooth supersingular
 states. Inequalities are not claimed to be equalities.}
 \label{main:state-orders}
\end{table}

\subsection{Relation to existing mechanisms and scope}

The initial-value surface and its eight-point geometry originate in
the q-Painlev\'e literature \cite{JoshiLobb2016,JoshiRoffelsen2026}.
Anticanonical-cycle periods and the relation between a normal-bundle
torsion order and a Halphen pencil are established tools
\cite{GrossHackingKeel2015,CantatDolgachev2012}. We need their values
on the original integer charts, and must separately exclude
quasi-elliptic behavior in the allowed small characteristics.
Spectral line-bundle correspondences and cyclic Picard constructions
are likewise classical \cite{Beauville1990,InoueVanhaeckeYamazaki2015}.
The proof here supplies the actual state inverse, excludes a hidden
inseparable degree or pullback kernel, and descends the torsor action
over the original field. Minimal-model uniqueness
\cite[Tags 0C6B and 0C9Z]{StacksProject2026} is applied only after
both models over the same local energy ring have been shown minimal.

Hasse powers and semilinear iteration are not new mechanisms.
Cartier conventions require care even for one-dimensional spaces
\cite{AchterHowe2019}. Vlasenko's higher Hasse--Witt congruences already
include the residue iteration, and the formal-group integrality
statement in the cited author version is not restricted to ordinary
fibers \cite[Theorems 1(i) and 2, author version v3]{Vlasenko2018}.
Our comparison uses that specified version; no theorem from it is a
proof premise below. Dwork-crystal and higher-Hasse constructions
give finer congruences under their respective invertibility and
Frobenius-lift hypotheses \cite{BeukersVlasenko2021,Vlasenko2024}.
We prove the needed coefficients and all block deformations directly,
including the supersingular locus and characteristic two.

The local division of Frobenius differences and the resulting
obstruction classes have a broad prior theory
\cite{DeligneIllusie1987,DupuyZureickBrown2019,DupuyKatzRabinoffZureickBrown2019}.
Recent Frobenius--Witt formulations also treat arithmetic
Kodaira--Spencer classes \cite{Shimada2026}. These theories do not by
themselves identify the specific $L_m$ section with the original trace
congruence used here. Nor is ramification alone a distinction from
that literature. Root-of-unity first nonzero ramified terms occur in
other specified systems, for example in quantum K-theoretic spectral
relations \cite{KoroteevSmirnov2026}; a spectral relation is not an
identity of the present two-direction critical ideals.
Lubin--Tate already give a characteristic-two height-two elliptic
lifting family \cite[Section 3.5]{LubinTate1966}. Accordingly, neither
height two, the number five, nor a normalizable unit parameter is
counted as a separate general mechanism here. The calculation concerns
the mixed ideal and the original base-parameter action in this model.

All topic-specific proofs needed for the three theorems are given in
the body. Sections~\ref{geom:section}--\ref{hasse:section} establish
the original geometric and Hasse identifications;
Sections~\ref{trace:section}--\ref{first:section} prove the integral
trace and complete first-layer ideal; Sections~\ref{odd:section}
and \ref{two:section} give the two higher branches. The paper does
not determine higher full thickness, singular-level critical ideals,
or exact orders after extra ramified state extensions. It makes no
general crystalline or formal-group comparison and no claim about
Riemann zeros.
